\documentclass[12pt]{article}
\usepackage[a4paper,margin=1in]{geometry}
\usepackage[utf8]{inputenc}
\usepackage[T1]{fontenc}
\usepackage{amsmath,amssymb}
\usepackage{textcomp}
\usepackage{parskip}
\begin{document}
\begin{center}
{\Large \textbf{2005 STEP II}}
\end{center}

\subsection*{Page 1}
\begin{flushleft}
STEP II, 2005 2\par
\bigskip\textbf{Section A: Pure Mathematics}\par
\bigskip\textbf{Question 1.} Find the three values of x for which the derivative of x\par
e\par
\[x^2\]
is zero.\par
-\par
Given that a and b are distinct positive numbers, find a polynomial P(x) such that the deriva-\par
tive of P(x)e\par
\[x^2\]
\[is zero for x = 0, x = a and x = b , but for no other values of x.\]
-\par
± ±\par
\bigskip\textbf{Question 2.} For any positive integer N , the function f (N ) is defined by\par
1 1 1\par
\[f (N ) = N 1 1 1\]
\[- p^1 - p^2 \cdots - p\]
k\par
\gamma \delta\gamma \delta \gamma \delta\par
\[where p^1, p^2, . . . , p\]
k\par
are the only prime numbers that are factors of N .\par
1 1\par
\[Thus f (80) = 80(1 )(1 ) .\]
2 5\par
- -\par
(a) (i) Evaluate f (12) and f (180).\par
\textbf{(ii)} Show that f (N ) is an integer for all N .\par
(b) Prove, or disprove by means of a counterexample, each of the following:\par
\textbf{(i)} f (m)f (n) = f (mn) ;\par
\textbf{(ii)} f (p)f (q) = f (pq) if p and q are distinct prime numbers;\par
\textbf{(iii)} f (p)f (q) = f (pq) only if p and q are distinct prime numbers.\par
\[(c) Find a positive integer m and a prime number p such that f (pm) = 146410 .\]
\[3 Give a sketch, for 0 6 x 6 \pi/2, of the curve\]
\[y = (sin x x cos x) ,\]
-\par
and show that 0 6 y 6 1 .\par
Show that:\par
\[\pi/2\]
\[\pi\]
\textbf{(i)} y dx = 2 ;\par
0 - 2\par
\[\int\]
\[\pi/2\]
\[\pi\]
\[\pi\]
\textbf{(ii)} y dx = .\par
0 48 - 8\par
\[\int\]
\[Deduce that \pi + 18\pi < 96 .\]
\end{flushleft}
\subsection*{Page 2}
\begin{flushleft}
STEP II, 2005 3\par
\bigskip\textbf{Question 4.} The positive numbers a, b and c satisfy bc = a + 1. Prove that\par
1 1 1\par
\[arctan + arctan = arctan .\]
a + b a + c a\par
\lambda \mu \lambda \mu \lambda \mu\par
The positive numbers p, q, r, s, t, u and v satisfy\par
2 2 2\par
\[st = (p + q) + 1 , uv = (p + r) + 1 , qr = p + 1 .\]
Prove that\par
1 1 1 1 1\par
\[arctan + arctan + arctan + arctan = arctan .\]
p + q + s p + q + t p + r + u p + r + v p\par
\lambda \mu \lambda \mu \lambda \mu \lambda \mu \lambda \mu\par
Hence show that\par
1 1 1 1 1\par
\[arctan + arctan + arctan + arctan = arctan .\]
13 21 82 187 7\par
\lambda \mu \lambda \mu \lambda \mu \lambda \mu \lambda \mu\par
[ Note that arctan x is another notation for tan x . ]\par
-\par
\bigskip\textbf{Question 5.} The angle A of triangle ABC is a right angle and the sides BC, CA and AB are of lengths a,\par
b and c, respectively. Each side of the triangle is tangent to the circle S1 which is of radius r.\par
\[Show that 2r = b + c a.\]
-\par
Each vertex of the triangle lies on the circle S2. The ratio of the area of the region between S1\par
\[and the triangle to the area of S^2 is denoted by R . Show that\]
\[\piR = (\pi 1)q + 2\piq (\pi + 1) ,\]
- - -\par
b + c\par
\[where q = . Deduce that\]
a\par
R 6 .\par
\[\pi(\pi 1)\]
-\par
1 2\par
6 (i) Write down the general term in the expansion in powers of x of (1 x) , (1 x) and\par
- -\par
3 - -\par
\[(1 x) , where x < 1 .\]
-\par
- | |\par
\[Evaluate \infty n^2 n and \infty n 2 2 n .\]
- -\par
\[n=1 n=1\]
X X\par
1 \infty (2n)! xn\par
\textbf{(ii)} Show that (1\par
-\par
\[x) - 2 = (n!)2 22n , for\]
|\par
x\par
|\par
\[< 1 .\]
\[n=0\]
X\par
\infty (2n)! \infty n(2n)!\par
Evaluate and .\par
\[(n!)222n3n (n!)222n3n\]
\[n=0 n=1\]
X X\par
\end{flushleft}
\subsection*{Page 3}
\begin{flushleft}
STEP II, 2005 4\par
\bigskip\textbf{Question 7.} The position vectors, relative to an origin O, at time t of the particles P and Q are\par
\[1 3 3\sqrt\{3\} 1\]
\[cos t i + sin t j + 0 k and cos(t + \pi) i + k + 3 sin(t + \pi) j ,\]
4 2 2 4\par
\[respectively, where 0 6 t 6 2\pi . \leq \geq\]
\textbf{(i)} Give a geometrical description of the motion of P and Q.\par
\textbf{(ii)} Let θ be the angle P OQ at time t that satisfies 0 6 θ 6 \pi . Show that\par
\[3\sqrt\{2\} 1 1\]
\[cos θ = cos(2t + \pi) .\]
8 4 4\par
-\par
\textbf{(iii)} Show that the total time for which θ > 1 \pi is 3 \pi .\par
4 2\par
\bigskip\textbf{Question 8.} For x > 0 the curve C is defined by\par
3 2\par
dy x y\par
\[=\]
\[dx (1 + x^2)5/2\]
\[with y = 1 when x = 0 . Show that\]
1 2 + 3x 1\par
\[= +\]
\[y 3(1 + x^2)3/2 3\]
and hence that for large positive x\par
y 3 .\par
\approx - x\par
Draw a sketch of C.\par
\[On a separate diagram draw a sketch of the two curves defined for x > 0 by\]
3 3\par
dz x z\par
\[=\]
\[dx 2(1 + x^2)5/2\]
\[with z = 1 at x = 0 on one curve, and z = 1 at x = 0 on the other.\]
-\par
\end{flushleft}
\subsection*{Page 4}
\begin{flushleft}
STEP II, 2005 5\par
\bigskip\textbf{Section B: Mechanics}\par
\bigskip\textbf{Question 9.} Two particles, A and B, of masses m and 2m, respectively, are placed on a line of greatest\par
slope, `, of a rough inclined plane which makes an angle of 30 with the horizontal. The\par
\textasciicircum{}\circ\par
coefficient of friction between A and the plane is 1 \sqrt3 and the coefficient of friction between B\par
and the plane is 1 \sqrt3. The particles are at rest with B higher up ` than A and are connected\par
by a light inextensible string which is taut. A force P is applied to B.\par
\textbf{(i)} Show that the least magnitude of P for which the two particles move upwards along `\par
\[is 11 \sqrt\{3\} mg and give, in this case, the direction in which P acts.\]
\textbf{(ii)} Find the least magnitude of P for which the particles do not slip downwards along `.\par
\bigskip\textbf{Question 10.} The points A and B are 180 metres apart and lie on horizontal ground. A missile is launched\par
1 3\par
from A at speed of 100 m s and at an acute angle of elevation to the line AB of arcsin . A\par
- 5\par
time T seconds later, an anti-missile missile is launched from B, at speed of 200 m s and at\par
-\par
an acute angle of elevation to the line BA of arcsin . The motion of both missiles takes place\par
in the vertical plane containing A and B, and the missiles collide.\par
\[Taking g = 10 m s and ignoring air resistance, find T .\]
-\par
3 1 3\par
[Note that arcsin is another notation for sin .]\par
5 - 5\par
\bigskip\textbf{Question 11.} A plane is inclined at an angle arctan to the horizontal and a small, smooth, light pulley P\par
is fixed to the top of the plane. A string, AP B, passes over the pulley. A particle of mass m1\par
is attached to the string at A and rests on the inclined plane with AP parallel to a line of\par
greatest slope in the plane. A particle of mass m2, where m2 > m1, is attached to the string\par
at B and hangs freely with BP vertical. The coefficient of friction between the particle at A\par
and the plane is .\par
The system is released from rest with the string taut. Show that the acceleration of the\par
\[m^2 m^1\]
particles is - g.\par
\[m^2 + m^1\]
At a time T after release, the string breaks. Given that the particle at A does not reach the\par
pulley at any point in its motion, find an expression in terms of T for the time after release\par
at which the particle at A reaches its maximum height. It is found that, regardless of when\par
the string broke, this time is equal to the time taken by the particle at A to descend from its\par
\[point of maximum height to the point at which it was released. Find the ratio m^1 : m^2.\]
3 1 3\par
[Note that arctan is another notation for tan .]\par
4 - 4\par
\end{flushleft}
\subsection*{Page 5}
\begin{flushleft}
STEP II, 2005 6\par
\bigskip\textbf{Section C: Probability and Statistics}\par
\bigskip\textbf{Question 12.} The twins Anna and Bella share a computer and never sign their e-mails. When I e-mail\par
them, only the twin currently online responds. The probability that it is Anna who is online\par
is p and she answers each question I ask her truthfully with probability a, independently of\par
all her other answers, even if a question is repeated. The probability that it is Bella who\par
is online is q, where q = 1 p, and she answers each question truthfully with probability b,\par
-\par
independently of all her other answers, even if a question is repeated.\par
\textbf{(i)} I send the twins the e-mail: 'Toss a fair coin and answer the following question. Did the\par
coin come down heads?'. I receive the answer 'yes'. Show that the probability that the\par
\[coin did come down heads is if and only if 2(ap + bq) = 1.\]
\textbf{(ii)} I send the twins the e-mail: 'Toss a fair coin and answer the following question. Did\par
the coin come down heads?'. I receive the answer 'yes'. I then send the e-mail: 'Did\par
the coin come down heads?' and I receive the answer 'no'. Show that the probability\par
(taking into account these answers) that the coin did come down heads is .\par
\textbf{(iii)} I send the twins the e-mail: 'Toss a fair coin and answer the following question. Did\par
the coin come down heads?'. I receive the answer 'yes'. I then send the e-mail: 'Did the\par
\[coin come down heads?' and I receive the answer 'yes'. Show that, if 2(ap + bq) = 1,\]
the probability (taking into account these answers) that the coin did come down heads\par
is .\par
\bigskip\textbf{Question 13.} The number of printing errors on any page of a large book of N pages is modelled by a\par
Poisson variate with parameter λ and is statistically independent of the number of printing\par
errors on any other page. The number of pages in a random sample of n pages (where n is\par
\[much smaller than N and n > 2) which contain fewer than two errors is denoted by Y . Show\]
\[that P(Y = k) = n pkqn k where p = (1 + λ)e λ and q = 1 p .\]
k - - -\par
Show also that, ifλis sufficiently small,\par
1 2\par
\textbf{(i)} q λ ;\par
\approx\par
\textbf{(ii)} the largest value of n for which P(Y = n) > 1 λ is approximately 2/λ ;\par
-\par
\textbf{(iii)} P(Y > 1 Y > 0) 1 n(λ 2 /2)n 1 .\par
-\par
| \approx -\par
\end{flushleft}
\subsection*{Page 6}
\begin{flushleft}
STEP II, 2005 7\par
\bigskip\textbf{Question 14.} The probability density function f (x) of the random variable X is given by\par
\[f (x) = k [φ(x) + λg (x)] ,\]
where φ(x) is the probability density function of a normal variate with mean 0 and variance 1,\par
λ is a positive constant, and g (x) is a probability density function defined by\par
1/λ for 0 6 x 6 λ ;\par
\[g (x) =\]
(0 otherwise.\par
Find µ, the mean of X, in terms of λ, and prove that σ, the standard deviation of X, satisfies.\par
4 3\par
λ + 4λ + 12λ + 12\par
\[σ = .\]
12(1 + λ)2\par
\[In the case λ = 2:\]
\textbf{(i)} draw a sketch of the curve y = f (x);\par
\textbf{(ii)} express the cumulative distribution function of X in terms of Φ(x), the cumulative\par
distribution function corresponding to φ(x);\par
2 2\par
\textbf{(iii)} evaluate P(0 < X < µ + 2σ), given that Φ( + \sqrt7) = 0.9921.\par
3 3\par
\end{flushleft}
\end{document}
